\id{МРНТИ 61.13.15}{}

\begin{header}
\swa{А.К.~Базанова, А.К.~Какимов, Н.К.~Ибрагимов, М.М.~Ташыбаева}{ИССЛЕДОВАНИЕ ПАРАМЕТРОВ РАБОТЫ СТЕНДА ДЛЯ АДСОРБЦИОННО-ФИЛЬТРАЦИОННОЙ ОЧИСТКИ ОТ ТОКСИЧНЫХ ВЕЩЕСТВ НА ОСНОВЕ ПРИРОДНЫХ СОРБЕНТОВ}

А.К.~Базанова\envelope,
А.К.~Какимов,
Н.К.~Ибрагимов,
М.М.~Ташыбаева
\end{header}

\begin{affil}
«Шәкәрім университеті», Семей, Казахстан

\corrauthor{Корреспондент-автор: arayka.bazanovak@mail.ru}
\end{affil}

В данной работе исследуется параметры функционирования экспериментальной
установки ад\-сорбционно-фильтрационной очистки жидких пищевых продуктов
от токсичных веществ с использованием природных сорбентов (цеолита,
шунгита и кокосового активированного угля). Установка представляет собой
напорную систему с тремя сорбционными колоннами, которые могут
включаться индивидуально или в различных комбинациях, что позволяет
варьировать режимы фильтрации. В ходе эксперимента изучалось влияние
объёма сорбентов (50, 100, 150 и 200 мл по шкале мерного стакана) и
мощности насоса (135, 200 и 245 Вт), а также самотёчного режима на
скорость процесса. Экспериментальные данные показали, что оптимальные
результаты достигаются при мощности насоса 200 Вт с использованием
цеолита (вторая колонна) и кокосового активированного угля (третья
колонна), обеспечивающих наименьшее время фильтрации при сохранении
высокой эффективности очистки. Наиболее эффективными сорбентами
оказались цеолит и кокосовый активированный уголь, отличающиеся развитой
пористой структурой и высокой удельной поверхностью. Самотёчный режим
показал наименьшую эффективность. Полученные результаты подтверждают
целесообразность применения природных сорбентов для повышения
безопасности жидких пищевых продуктов и могут быть использованы при
проектировании промышленных систем предварительной очистки.

{\bfseries Ключевые слова:} адсорбционно-фильтрационный стенд, фильтрация,
адсорбция, природные сорбенты, токсичные вещества, пищевая безопасность.

\begin{header}
STUDY OF THE OPERATING PARAMETERS OF THE STAND FOR ADSORPTION-FILTRATION PURIFICATION FROM TOXIC SUBSTANCES BASED ON NATURAL SORBENTS

A.K.~Bazanova\envelope,
A.K.~Kakimov,
N.K.~Ibragimov,
M.M.~Tashybayeva
\end{header}

\begin{affil}
«Shakarim University», Semey, Kazakhstan,

e-mail: arayka.bazanovak@mail.ru
\end{affil}

This article investigates the operating parameters of an experimental
adsorption-filtration unit for the purification of liquid food products
from toxic substances using natural sorbents (zeolite, shungite, and
coconut-based activated carbon). The unit is designed as a pressure
system with three sorption columns that can be operated individually or
in various combinations, allowing for flexible adjustment of filtration
modes. The experiments examined the effect of sorbent volume (50, 100,
150, and 200 ml according to the measuring scale) and pump power (135,
200, and 245 W), as well as gravity flow mode, on the filtration rate.
The results demonstrated that the most efficient performance was
achieved at a pump power of 200 W with the use of zeolite (second
column) and coconut activated carbon (third column), providing the
shortest filtration time while maintaining high purification efficiency.
Zeolite and coconut activated carbon proved to be the most effective
sorbents due to their developed porous structure and high specific
surface area. The gravity flow mode showed the lowest efficiency. The
findings confirm the feasibility of applying natural sorbents to improve
the safety of liquid food products and their potential use in the design
of industrial systems for preliminary purification.

{\bfseries Keywords:} adsorption-filtration unit, filtration, adsorption,
natural sorbents, toxic substances, food safety.

\begin{header}
ТАБИҒИ СОРБЕНТТЕР НЕГІЗІНДЕ УЛЫ ЗАТТАРДАН АДСОРБЦИЯЛЫҚ-ФИЛЬТРАЦИЯЛЫҚ ТАЗАРТУ СТЕНДІ ЖҰМЫСЫНЫҢ ПАРАМЕТРЛЕРІН ЗЕРТТЕУ

А.Қ.~Базанова\envelope,
А.К.~Кәкімов,
Н.К.~Ибрагимов,
М.М.~Ташыбаева
\end{header}

\begin{affil}
«Шәкәрім университеті», Семей, Қазақстан,

e-mail: arayka.bazanovak@mail.ru
\end{affil}

Бұл зерттеу жұмысында табиғи сорбенттерді (цеолит, шунгит және кокос
қабығынан алынған белсендірілген көмір) қолдана отырып, сұйық тағам
өнімдерінен улы (зиянды) заттарды адсорб\-циялық-фильтрациялық тазартуға
арналған эксперименттік қондырғының жұмыс параметрлері зерттелді.
Қондырғы үш сорбциялық бағандардан тұратын қысымдық жүйе, оларды жеке
немесе түрлі комбинацияларда қосу арқылы фильтрация режимдерін өзгертуге
мүмкіндік береді. Эксперимент жүргізу барысында сорбенттердің көлемі
(50, 100, 150 және 200 мл өлшеуіш стақан шкаласы бойынша) және сорғы
қуаты (135, 200 және 245 Вт), сондай-ақ өздігінен ағу жұмыс режимінің
процестің жылдамдығына әсері зерттелді. Эксперименттік көрсеткіштер
негізінде, оңтайлы нәтижелер 200 Вт сорғы қуатында, цеолитті (екінші
баған) және кокос қабығынан алынған белсендірілген көмірін (үшінші
баған) қолданғанда алынды, бұл жағдайда фильтрация уақыты ең төмен
болып, тазарту тиімділігі жоғары деңгейде сақталды. Ең тиімді сорбенттер
цеолит пен кокос қабығынан алынған белсендірілген көмір болып шықты,
олар кеуекті құрылымымен және жоғары меншікті бетімен ерекшеленеді.
Өздігінен ағу жұмыс режимі ең төмен тиімділікті көрсетті. Алынған
нәтижелер сұйық тағам өнімдерінің қауіпсіздігін арттыру үшін табиғи
сорбенттерді қолданудың мақсатқа сай екенін дәлелдейді және өнеркәсіптік
алдын-ала тазарту жүйелерін жобалауда пайдаланылуы мүмкін.

{\bfseries Түйін сөздер:} адсорбциялық-фильтрациялық стенд, фильтрация,
адсорбция, табиғи сорбенттер, улы (зиянды) заттар, тағам қауіпсіздігі.

\begin{multicols}{2}
{\bfseries Введение.} Эффективным подходом к снижению токсичных веществ в
жидких пищевых продуктах являются адсорбционно-фильтрационные процессы.
В связи с этим актуальным представляется использование природных
сорбентов, отличающихся меньшей стоимостью по сравнению с коммерческими
аналогами, а также высокой эффективностью и экологической безопасностью
{[}1{]}. Фильтрованием называют процесс разделения суспензий с
использованием пористых перегородок, которые задерживают твёрдую фазу
суспензии и пропускают её жидкую фазу. Разделение суспензии, состоящей
из жидкости и взвешенных в ней твёрдых частиц, производят при помощи
фильтра, который в простейшем виде является сосудом, разделённым на две
части фильтровальной перегородкой. Процесс адсорбции заключается в том,
что твёрдое вещество (адсорбент) удерживает на своей поверхности один
или несколько компонентов из газа или раствора {[}2{]}. В результате на
поверхности адсорбента образуется слой адсорбата. Притяжение обусловлено
как силами Ван-дер-Ваальса, так и взаимодействиями, основанными на
гидрофобном отталкивании {[}3{]}. Адсорбция является эффективным методом
удаления токсичных веществ из жидких пищевых продуктах. Метод отличается
рядом преимуществ по сравнению с альтернативными технологиями очистки:

- низкая стоимость строительства и эксплуатации очистных систем;

- высокая эффективность при удалении концентрированных загрязнителей;

- возможность регенерации сорбента и повторного использования;

- способность к адсорбции многокомпонентных смесей веществ {[}4{]}.

Радиационная опасность во многом связана с загрязнением питьевой воды и
пищевых продуктов радионуклидами и тяжёлыми металлами, что представляет
собой актуальную проблему пищевой безопасности {[}5{]}. Цель
работы-исследовать параметры функционирования
адсорбционно-фильтрационной установки с применением природных сорбентов.
Экспериментальная установка, представленная в исследовании, обеспечивает
эффективную адсорбционную очистку жидких пищевых продуктов и
способствует повышению их безопасности.

{\bfseries Материалы и методы.} С целью исследования адсорбционных свойств
была использована экспериментальная установка, сконструированная на
кафедре «Биоинженерных систем» университета Шәкәрім. Наименование данной
установки является стенд {[}6{]} для фильтрационной очистки жидких
пищевых продуктов от радионуклидов и солей тяжёлых металлов. Она
представляет собой напорную систему с тремя вертикальными колоннами,
соединёнными с насосом, трубопроводом и системой управляющими кранами
{[}7{]}.
\end{multicols}

\fig[0.45\textwidth]{p/image35}[Рис.1 - Экспериментальная установка для адсорбционной очистки
жидких пищевых продуктов от токсичных веществ {[}7{]}]

\begin{multicols}{2}
1-опорная конструкция, 2-приёмная ёмкость, 3-кран для подачи жидкости
самотёком, 4-кран для подачи жидкости с помощью насоса (под давлением),
5-центробежно-лопастной насос, 6-блок питания, 7, 8, 9-сорбционные
колонны, 10-17-управляющие краны, 18-23-прокладки из технической резины,
устойчивые к высоким/низким температурам, кислотам и щелочам, 24, 25,
26-сетки, предотвращающие вымывание сорбента, 27, 28-накопительные
ёмкости, 29-манометр, 30-расходомер.

Преимуществом данной установки является возможность как индивидуального
включения каждой колонны, так и их комбинированного соединения в
различных последовательностях, что обеспечивает гибкость эксперимента и
позволяет подбирать наиболее эффективные режимы адсорбции. В
исследовании были опробованы разные виды природных сорбентов, что дало
возможность определить наиболее подходящие материалы для снижения
содержание токсичных веществ. В качестве сорбентов использовались:
цеолит {[}8{]} Чанканайского месторождения (Кербулакский район,
Алматинская область), шунгит {[}9{]} Коксуского месторождения и
активированный уголь марки 207С, полученный из скорлупы кокосовых орехов
{[}10{]}. На рисунке 2 представлены сорбенты, применённые в
эксперименте. Каждая колонна загружалась определённым видом сорбента
(фракция 2-1 мм), отобранным с учётом его селективности к радионуклидам
и ионам тяжёлых металлов.
\end{multicols}

\begin{figs}[Рис.2 - Чанканайский цеолит № 1, коксуский шунгит № 2,
активированный уголь из скорлупы кокосовых орехов № 3 {[}8-10{]}]
    \fig[0.32\textwidth][\textwidth]{p/image36}[1]
    \fig[0.32\textwidth][\textwidth]{p/image37}[2]
    \fig[0.32\textwidth][\textwidth]{p/image38}[3]
\end{figs}

\begin{multicols}{2}
\emph{Объекты исследования}. В качестве объекта исследования
была выбрана питьевая вода из водопроводной сети. Для оценки параметров
работы экспериментального стенда использовали 5 литров воды при средней
температуре 18-20 °С. В ходе эксперимента применялись природные сорбенты
различного объёма -50, 100, 150 и 200 мл (по шкале мерного стакана).
Длительность процесса работы адсорбционно-фильтрационной установки
составляла 3-5 минут.

{\bfseries Результаты и обсуждение.} В первой колонне был загружен шунгит,
во второй-цеолит, в третьей-кокосовый активированный уголь. Насос имел
три уровня рабочей мощности: первая мощность-135 Вт, вторая мощность-200
Вт, третья мощность-245 Вт.

На графиках (рисунки 3-6) по оси X представлен объём сорбентов (50, 100,
150 и 200 мл по шкале мерного стакана), а по оси Y-время процесса (сек).
Каждое исследование отражено отдельной диаграммой с указанием
соответствующего режима фильтрации.

Экспериментальные данные показали (рисунок 3), что время фильтрации
возрастает с увеличением объёма сорбентов. Наиболее эффективными
оказались цеолит (2 колонна), кокосовый активированный уголь (3 колонна)
объясняется их развитой пористой структурой и высокой удельной
поверхностью, а также их совместное применение (1+2+3 колонны). Однако
при объёме выше 150 мл наблюдалось резкое увеличение времени процесса,
особенно при комбинированном использовании трёх сорбентов.

При увеличении мощности до 200 Вт общее время фильтрации (рисунок 4)
сократилось по сравнению с режимом 135 Вт. Оптимальные результаты были
достигнуты при использовании цеолита и кокосового активированного угля
(2 и 3 колонны), а также при комбинированной загрузке (1+2+3 колонны).
Рост времени при увеличении объёма был менее выраженным, что
свидетельствует о более стабильном протекании процесса.
\end{multicols}

\fig[0.65\textwidth]{p/image39}[Рис.3 - Фильтрация воды при мощности насоса 135 Вт (первая
скорость) с использованием различных сорбентов]

\fig[0.65\textwidth]{p/image40}[Рис.4 - Фильтрация воды при мощности насоса 200 Вт (вторая
скорость) с использованием различных сорбентов]

\fig[0.65\textwidth]{p/image41}[Рис.5 - Фильтрация воды при мощности насоса 245 Вт (третья
скорость) с использованием различных сорбентов]

\begin{multicols}{2}
В данном режиме процесс проходил намного быстрее: значения времени
фильтрации (рисунок 5) оказались минимальными среди насосных режимов.
Особенно эффективно проявил себя кокосовый активированный уголь (3
колонна) и его комбинация с цеолитом (2+3 колонны). Однако при больших
объёмах сорбентов наблюдалась тенденция к перегрузке и росту времени
процесса.
\end{multicols}

\fig[0.65\textwidth]{p/image42}[Рис.6 - Фильтрация воды в режиме самотёка (без насоса) с
использованием различных сорбентов]

\begin{multicols}{2}
Данный вариант на рисунке 6 показал наихудшие результаты: время
фильтрации оказалось в десятки раз выше по сравнению с насосными
режимами и достигало 300-370 с при объёме 200 мл. Наиболее выраженное
замедление наблюдалось при комбинации шунгита и цеолита (1+2 колонны), а
также при использовании всех трёх сорбентов одновременно (1+2+3
колонны). Таким образом, самотёчный режим можно признать неприемлемым
для практического применения.

{\bfseries Выводы.} Проведён анализ и определены, что:

- самотёчный режим неэффективен из-за чрезмерного увеличения времени
фильтрации.

- при мощности 135 Вт (первая скорость) наблюдается рост времени при
увеличении

объёма, что ограничивает применение режима;

- наиболее сбалансированные результаты получены при мощности \emph{200
Вт (2}

\emph{скорость):} процесс идёт значительно быстрее, чем при 135 Вт
(первая скорость), но без резкого увеличения времени, как при 245 Вт
(третья скорость);

- при 245 Вт (третья скорость) достигается минимальное время фильтрации,
но при

больших объёмах сорбентов появляется перегрузка и резкий рост времени.

Следовательно, \emph{оптимальным режимом фильтрации для установки
является работа при мощности 200 Вт (2 скорость),} с применением
цеолита, кокосового активированного угля (2+3 колонны) или их комбинации
(1+2+3 колонны). Наиболее стабильные результаты получены при объёмах
сорбентов \emph{100-150 мл} по шкале мерного стакана, где обеспечивается
минимальное время фильтрации при сохранении высокой эффективности
очистки. При объёме 200 мл наблюдалась тенденция к перегрузке, что
ограничивает практическое применение больших загрузок. Полученные
результаты позволяют рекомендовать усовершенствованную установку для
применения в системах предварительной очистки жидких пищевых продуктов
от токсичных веществ.
\end{multicols}

\begin{center}
{\bfseries Литература}
\end{center}

\begin{refs}
1. Dolenko S.O., Trifonova M.Y., Ivanova Z.G. Influence of Natural Water
Components on the Sorption of Dodecyl Sulfate Anions on Kaolinite//J.
Water Chem. Technol.-2022.-Vol.44(1).-P.1-9. DOI
10.3103/S1063455X22010027.

2. Мантлер С.Н., Жуманазарова Г.М. Процессы и аппараты химической
технологии: Учебное пособие.Алматы: «Бастау».-2018.-Р.256. ISBN
978-601-7275-76-1.

3. Chen Y. et al. Application studies of activated carbon derived from
rice husks produced by chemical-thermal process - A review//Advances
in colloid and interface science. - 2011. -Vol.163(1).- P.39-52. DOI
\href{https://doi.org/10.1016/j.cis.2011.01.006}{10.1016/j.cis.2011.01.006}.

4. Беляева Е.А. Адсорбционная очистка сточных вод//Синтез науки и
образования в решении глобальных проблем современности: сборник статей
Международной научно-практической
конференции//Стерлитамак:АМИ.-2020.-С.57-58.-URL:
\url{https://ami.im/sbornik/MNPK-272.pdf}.

5. Kakimov A., Jilkisheva A., Mayorov A., Zharykbasova K., Kakimova Z.,
Mirasheva G., Zharykbasov Y., Zolotov A.,Yessimbekov Z. Developing the
process parameters of milk pasteurization for reducing the
concentration of toxic elements and radionuclides//ARPN Journal of
Engineering and Applied Sciences.-2019.Vol.14(13).-P.2443-2447.

6. Стенд для фильтрационной очистки жидких пищевых продуктов от
радионуклидов и солей тяжёлых металлов. Патент РК №11024 G01N 15/00.

7. Базанова А.К., Какимов А.К., Лобасенко Б.А. и другие.
Совершенствование установки для снижения содержания токсичных веществ
в молоке// Вестник АТУ -2025.-Т.149(3).-С.19-26.
\href{https://doi.org/10.48184/2304-568X-2025-3-19-26}{DOI
10.48184/2304-568X-2025-3-19-26}.

8. Yurak V., Apakashev R., Dushin A., Usmanov A., Lebzin M., Malyshev A.
Testing of Natural Sorbents for the Assessment of Heavy Metal Ions'
Adsorption.//Appl. Sci. -2021.-Vol 11(8):3723. DOI
10.3390/app11083723.

9. Мусина У.Ш. Экологический потенциал Коксуского
шунгита//Гидрометеорология и экология.- 2010.- № 4.- С.154-159.

10. Кокосовый активированный уголь (сорбция):-URL:
\url{https://sebekpro.ru/katalog-resheniy/gotovye-resheniya/172-kokosovyj-aktivirovannyj-ugol-sorbtsiya.html.-}
Дата обращения: 29.09.2025.
\end{refs}

\begin{center}
{\bfseries References}
\end{center}

\begin{refs}
1. Dolenko S.O., Trifonova M.Y., Ivanova, Z.G. Influence of Natural
Water Components on the Sorption of Dodecyl Sulfate Anions on
Kaolinite//J. Water Chem. Technol.-2022.-Vol.44(1). -P.1-9. DOI
10.3103/S1063455X22010027.

2. Mantler S.N., Zhumanazarova G.M. Processy i apparaty himicheskoj
tehnologii: Uchebnoe posobie. Almaty: «Bastau».-2018.-R.256. ISBN
978-601-7275-76-1. {[}in Russian{]}

3. Chen Y. et al. Application studies of activated carbon derived from
rice husks produced by chemical-thermal process - A review//Advances in
colloid and interface science. - 2011. -Vol.163(1).-P.39-52. DOI
\href{https://doi.org/10.1016/j.cis.2011.01.006}{10.1016/j.cis.2011.01.006}.

4. Beljaeva E.A. Adsorbcionnaja ochistka stochnyh vod//Sintez nauki i
obrazovanija v reshenii global' nyh problem
sovremennosti: sbornik statej Mezhdunarodnoj nauchno-prakticheskoj
konferencii// Sterlitamak: AMI.-2020.-S.57-58.-URL:
https://ami.im/sbornik/MNPK-272.pdf. {[}in Russian{]}

5. Kakimov A., Jilkisheva A., Mayorov A., Zharykbasova K., Kakimova Z.,
Mirasheva G., Zharykbasov Y., Zolotov A.,Yessimbekov Z. Developing the
process parameters of milk pasteurization for reducing the concentration
of toxic elements and radionuclides//ARPN Journal of Engineering and
Applied Sciences.-2019.Vol.14(13).-P.2443-2447.

6. Stend dlja fil' tracionnoj ochistki zhidkih pishhevyh
produktov ot radionuklidov i solej tjazhjolyh metallov. Patent RK №11024
G01N 15/00. {[}in Russian{]}

7. Bazanova A.K., Kakimov A.K., Lobasenko B.A. i drugie.
Sovershenstvovanie ustanovki dlja snizhenija soderzhanija toksichnyh
veshhestv v moloke// Vestnik ATU -2025.-T.149(3).-S.19-26. DOI
10.48184/2304-568X-2025-3-19-26. {[}in Russian{]}

8. Yurak V., Apakashev R., Dushin A., Usmanov A., Lebzin M., Malyshev A.
Testing of Natural Sorbents for the Assessment of Heavy Metal Ions'
Adsorption.//Appl. Sci. -2021.-Vol 11(8):3723. DOI 10.3390/app11083723.

9. Musina U.Sh. Jekologicheskij potencial Koksuskogo
shungita//Gidrometeorologija i jekologija.- 2010.- № 4.- S.154-159.
{[}in Russian{]}

10. Kokosovyj aktivirovannyj ugol'{} (sorbcija):- URL:
https://sebekpro.ru/katalog-resheniy/gotovye-resheniya/172-kokosovyj-aktivirovannyj-ugol-sorbtsiya.html.-
Data obrashhenija: 29.09.2025. {[}in Russian{]}
\end{refs}

\begin{info}
\hspace{1em}\emph{{\bfseries Сведения об авторах}}

Какимов А.К.-доктор технических наук, профессор, Шәкәрім университет,
Семей, Казахстан, е-mail: bibi.53@mail.ru;

Ибрагимов Н.К.-кандидат технических наук, ассоциированный профессор,
преподаватель, Шәкәрім университет, Семей, Казахстан, e-mail:
ibragimnk@mail.ru;

Ташыбаева М.М.--PhD, преподаватель, Шәкәрім университет, Семей,
Казахстан, е-mail: marzhan06081990@gmail.com;

Базанова А.К.-докторант, Шәкәрім университет, Семей, Казахстан, e-mail:
arayka.bazanovak@mail.ru.

\hspace{1em}\emph{{\bfseries Information about authors}}

Kakimov A.K.-Doctor of Technical Sciences, Professor, Shakarim
University, Semey, Kazakhstan, e-mail: bibi.53@mail.ru;

Ibragimov N.K.-Candidate of Technical Sciences, Associate Professor,
Lecturer, Shakarim University, Semey, Kazakhstan, e-mail:
ibragimnk@mail.ru;

Tashybayeva M.M.-PhD, Lecturer, Shakarim University, Semey, Kazakhstan,
e-mail: marzhan06081990@gmail.com;

Bazanova A.K.- doctoral student, Shakarim University, Semey, Kazakhstan,
e-mail: arayka.bazanovak@mail.ru.
\end{info}
